%%%
%
% $Autor: Wings $
% $Datum: 2021-05-14 $
% $Pfad: GitLab/MLEdgeComputer $
% $Dateiname: 
% $Version: 4620 $
%
% !TeX spellcheck = de_DE/GB
% !TeX program = pdflatex
% !BIB program = biber/bibtex
% !TeX encoding = utf8
%
%%%

% source: https://www.digikey.fi/en/maker/projects/intro-to-tinyml-part-1-training-a-model-for-arduino-in-tensorflow/8f1fc8c0b83d417ab521c48864d2a8ec
% source: https://www.digikey.de/en/maker/projects/intro-to-tinyml-part-2-deploying-a-tensorflow-lite-model-to-arduino/59bf2d67256f4b40900a3fa670c14330

% Python-Code: E:\GitLab\MachineLearning\MLEdgeComputer\Code\TinyML\Short\SineWave.py
% Arduino-Code: E:\GitLab\MachineLearning\MLEdgeComputer\Code\TinyML\Short\SineWave.ino

\chapter{Introduction To TinyML}

\section{Training a Model for Arduino in TensorFlow}



When most of us think about artificial intelligence (AI) and machine learning (ML), we usually conjure up thoughts about home assistants, autonomous robots, and self-driving cars. However, the world of ML encompasses much more than that. Any time we train a mathematical model to help a computer make predictions, decisions, or classifications without being explicitly told to do so, we are using ML. 

In most cases, useful machine learning algorithms require copious amounts of computing resources (CPU cycles and RAM). However, TensorFlow Lite has recently released an experimental version that runs on several microcontrollers. Assuming we can create a model that fits on a device with limited resources, we can start turning embedded systems into tiny machine learning (TinyML) devices.

In this tutorial, we will create a neural network that is capable of predicting the output of the sine function. We will then convert this model to a TensorFlow Lite (TFLite) model and examine it using Netron.

If you would like to see an explanation of these steps in video form, check out this \href{https://www.youtube.com/watch?v=BzzqYNYOcWc&feature=emb_imp_woyt}{YouTube video}.

\subsection{Model Description}


We are going to create a 3-layer, fully connected neural network that predicts the output of the sine function. As such, we are using it as a regression model. 

\begin{figure}
  \includegraphics[width=\textwidth]{TinyML/Short/NN}
  \caption{Neural network model}
\end{figure}


The idea is to train a model that accepts values between $0$ and $2\pi$ and then outputs a value between $-1$ and $1$. If we label the input value as $x$ and the output as $y$, then the model should be able to predict $y = \sin(x)$.

\begin{figure}
  \includegraphics[width=\textwidth]{TinyML/Short/TrainingSteps}
  \caption{Neural network training steps}
\end{figure}



This might be one of the most inefficient, most roundabout ways to calculate a sinewave. However, it allows us to play with a small neural network with some nonlinearity and load it onto a microcontroller.

TensorFlow includes a converter class that allows us to convert a Keras model to a TensorFlow Lite model. The TensorFlow Lite model is stored as a \HREF{https://google.github.io/flatbuffers/}{FlatBuffer}, which is useful for reading large chunks of data one piece at a time (rather than having to load everything into RAM).

Once we have created and trained the model, we will run the TensorFlow Lite converter to create a tflite model. From there, we will need to store the tflite model as a series of bytes in a C constant array (either in a .c or .h file). This will allow us to load the model and use it for inference with the \HREF{https://www.tensorflow.org/lite/microcontrollers}{TensorFlow Lite for Microcontrollers} library.

\begin{figure}
  \includegraphics[width=\textwidth]{TinyML/Short/TensorFlowLite}
  \caption{TensorFlow Lite training and conversion}
\end{figure}


TensorFlow Lite has support for a few microcontroller boards, which you can see the list \href{https://www.tensorflow.org/lite/microcontrollers#supported_platforms}{here}. At the time this tutorial was released, only 8 microcontroller boards were supported. While there is an arduino library available for TensorFlow Lite, only the Nano 33 BLE Sense was supported. As a result, we will be using the \href{https://www.digikey.fi/en/products/detail/arduino/ABX00031/10239970?s=N4IgTCBcDaIIwAYCsCC0BBAQgDQXgzHKgHIAiIAugL5A}{Nano 33 BLE Sense} in the second part of this tutorial series.

\subsection{Google Colab}

Google Colab is a web-based Jupyter Notebook interface that runs on a Linux virtual machine (VM). It comes pre-loaded with most of the popular machine learning packages, and it's free!

Sign up for a \HREF{https://mail.google.com/mail/u/}{Google GMail account} and head to \URL{https://colab.research.google.com/}.

While you can play around with various machine learning algorithms in Colab, you're limited in VM resources and time. If you don’t engage with the interface for 90 min, you will be disconnected, and every 12 hours, the runtime will reset.

If you wish to train models for longer or need more resources, you will need to pay for the Pro version or set up your own machine learning development machine (local or server).

Note that most of the code below should also run locally in Python or on Jupyter Notebook. If you wish to run TensorFlow locally, please see this \HREF{https://www.digikey.fi/en/maker/projects/getting-started-with-machine-learning-using-tensorflow-and-keras/0746640deea84313998f5f95c8206e5b/}{tutorial}.

\subsection{Model Training}

Please note that the original code is based on Pete Warden’s work for TensorFlow Lite. I have tweaked the example some to work better on video.

Head to Colab and click New Notebook. Give that notebook a unique name, like ``tflite-sinewave-training.ipynb''. 

The code is best viewed in Notebook form so you can see example outputs. I will discuss some of the important pieces of code in this tutorial, but if you’re interested in running the Notebook yourself and dissecting the code, please see this Gist: \url{https://gist.github.com/ShawnHymel/79237fe6aee5a3653c497d879f746c0c}

At the top, we select the version of TensorFlow we wish to use. While I originally specified ``2.1'', this magic command only deals with major versions. As a result, it’s recommended to use:

\medskip

\textbf{\medskip \textbf{Copy Code}}
\PYTHON{\%tensorflow\_version 2.x}

\medskip

After that, we import our necessary packages and print their versions, which can be helpful when asking for help on various forums, as people will usually request to see which version of things you’re using.

We then specify various settings, including the number of samples we wish to use and the percentage of samples we wish to set aside for validation and test sets.

From there, we generate a random sampling of $x$ values to be used to train our model:



\medskip

 \textbf{\medskip \textbf{Copy Code}}
 
\begin{lstlisting}
# Generate some random samples
np.random.seed(1234)
x_values = np.random.uniform(low=0, high=(2 * math.pi), size=nsamples)
plt.plot(x_values)
We then calculate sin(x) for each sample and add a little bit of random, Gaussian noise to the output. This helps us ensure the model is an estimation and not a precise representation of a sinewave.
\end{lstlisting}


\medskip

 \textbf{\medskip \textbf{Copy Code}}
 
\begin{lstlisting}
# Create a noisy sinewave with these values
y_values = np.sin(x_values)   (0.1 * np.random.randn(x_values.shape[0]))
plt.plot(x_values, y_values, '.')
Next, we split the data up into training, validation, and test sets. It's usually a good idea to view (or plot) your input data and labels, if possible. Because input and output are only one value here (as opposed to a many-dimensional array), we can visualize the relationship between inputs and outputs more easily:
\end{lstlisting}



\medskip

 \textbf{\medskip \textbf{Copy Code}}
 
 
\begin{lstlisting}
# Split the dataset into training, validation, and test sets
val_split = int(val_ratio * nsamples)
test_split = int(val_split   (test_ratio * nsamples))
x_val, x_test, x_train = np.split(x_values, [val_split, test_split])
y_val, y_test, y_train = np.split(y_values, [val_split, test_split])

# Check that our splits add up correctly
assert(x_train.size   x_val.size   x_test.size) == nsamples

# Plot the data in each partition in different colors:
plt.plot(x_train, y_train, 'b.', label="Train")
plt.plot(x_test, y_test, 'r.', label="Test")
plt.plot(x_val, y_val, 'y.', label="Validate")
plt.legend()
plt.show()
\end{lstlisting}

This should give us a plot as shown here:

\begin{figure}
  \includegraphics[width=\textwidth]{TinyML/Short/SinePoints}
  \caption{Sinewave sets}
\end{figure}



Now that we have our data prepared, it’s time to create the model. We’ll use a 3-layer, fully-connected (dense) neural network.



\medskip

 \textbf{\medskip \textbf{Copy Code}}


\begin{lstlisting}
# Create a model
model = tf.keras.Sequential()
model.add(layers.Dense(16, activation='relu', input_shape=(1,)))
model.add(layers.Dense(16, activation='relu'))
model.add(layers.Dense(1))
\end{lstlisting}

Then, we can compile it with our set optimizers and loss functions, followed by the actual training procedure:



\medskip

 \textbf{\medskip \textbf{Copy Code}}


\begin{lstlisting}
# Add optimizer, loss function, and metrics to model and compile it
model.compile(optimizer='rmsprop', loss='mae', metrics=['mae'])
\end{lstlisting}


\medskip

 \textbf{\medskip \textbf{Copy Code}}


\begin{lstlisting}
# Train model
history = model.fit(x_train,
y_train,
epochs=500,
batch_size=100,
validation_data=(x_val, y_val))
\end{lstlisting}

After a few minutes, you should have a fully trained model. We can plot the model’s predicted sine function versus the test data to ensure that we have something that approximates a sinewave:



\medskip

 \textbf{\medskip \textbf{Copy Code}}
 
\begin{lstlisting}

# Plot the training history
loss = history.history['loss']
val_loss = history.history['val_loss']

epochs = range(1, len(loss)   1)

plt.plot(epochs, loss, 'bo', label='Training loss')
plt.plot(epochs, val_loss, 'b', label='Validation loss')
plt.title('Training and validation loss')
plt.legend()
plt.show()
\end{lstlisting}

The output of this code should give us a plot that looks like the following:

\begin{figure}
  \includegraphics[width=\textwidth]{TinyML/Short/SineWaveOutput}
  \caption{Sinewave model output}
\end{figure}


As you can see, our model should give us something that approximates a sinewave, even if it’s not exactly correct.

\subsection{Convert to TensorFlow Lite Model}

Now that we have a Keras model that has been trained, we need to convert it to something that our microcontroller can use. TensorFlow has a built-in converter function that will save the model as a TensorFlow Lite model file (stored as a FlatBuffer):



\medskip

\textbf{\medskip \textbf{Copy Code}}

\begin{lstlisting}
# Convert Keras model to a tflite model
converter = tf.lite.TFLiteConverter.from_keras_model(model)
converter.optimizations = [tf.lite.Optimize.OPTIMIZE_FOR_SIZE]
tflite_model = converter.convert()

open(tflite_model_name   '.tflite', 'wb').write(tflite_model)
\end{lstlisting}




Some advanced techniques, such as quantization, are available, which you can read about here: \url{https://www.tensorflow.org/lite/microcontrollers/build_convert}. Note that many of these are changing on a daily basis and may require a good deal of work to get running properly on a microcontroller. 

That being said, basic conversion seems to work for now, even if it does not offer the same level of memory saving that quantization offers.

To get our tflite file to run on a microcontroller, we need to save it as a constant byte array in a .h or .c file.

Note: for good C code, you should declare the byte array in a .h file and put the definition (the actual bytes) in a .c file. Because I want to keep this example simple, I will save everything in just a .h file.

This function can be used to read a tflite model file, convert it to hex bytes, and generate a .h file string that contains our model and any necessary C code:

\medskip

\textbf{\medskip \textbf{Copy Code}}

\begin{lstlisting}
# Function: Convert some hex value into an array for C programming
def hex_to_c_array(hex_data, var_name):

c_str = ''

# Create header guard
c_str  = '#ifndef '   var_name.upper()   '_H\n'
c_str  = '#define '   var_name.upper()   '_H\n\n'

# Add array length at top of file
c_str  = '\nunsigned int '   var_name   '_len = '   str(len(hex_data))   ';\n'

# Declare C variable
c_str  = 'unsigned char '   var_name   '[] = {'
    hex_array = []
    for i, val in enumerate(hex_data) :
    
    # Construct string from hex
    hex_str = format(val, '#04x')
    
    # Add formatting so each line stays within 80 characters
    if (i   1) < len(hex_data):
    hex_str  = ','
    if (i   1) % 12 == 0:
    hex_str  = '\n '
    hex_array.append(hex_str)
    
    # Add closing brace
    c_str  = '\n '   format(' '.join(hex_array))   '\n};\n\n'

# Close out header guard
c_str  = '#endif //'   var_name.upper()   '_H'

return c_str
\end{lstlisting}


We will call this function and save the \FILE{.h} file to our Colab VM:



\medskip

 \textbf{\medskip \textbf{Copy Code}}
 
\begin{lstlisting}
# Write TFLite model to a C source (or header) file
with open(c_model_name   '.h', 'w') as file:
file.write(hex_to_c_array(tflite_model, c_model_name))
\end{lstlisting}


\subsection{Save and Inspect Model File}

Click on the Files button in the left pane of the Colab window to inspect the files in the VM. You should see \FILE{sine\_model.h}. If not, click the Refresh button. Right-click on \FILE{sine\_model.h} and click Download.

\begin{figure}
  \includegraphics[width=\textwidth]{TinyML/Short/DownloadModelFile}
  \caption{Download model file}
\end{figure}



You can inspect the .h file to see that it contains legitimate C code (an array length and array of lots of hex values).

\begin{figure}
  \includegraphics[width=\textwidth]{TinyML/Short/InspectCModelFile}
  \caption{Inspect C model file}
\end{figure}


You can also download the \PYTHON{sine\_model.tflite} model file. If you do, I recommend installing \href{https://github.com/lutzroeder/netron}{Netron} to visualize the model. Simply open the .tflite file with Netron, and you should see a model description diagram. If you click on a layer, you can get information about that layer, such as the input and output types and tensor sizes.

\begin{figure}
  \includegraphics[width=\textwidth]{TinyML/Short/InspectModelWithNetron}
  \caption{Inspect model with Netron}
\end{figure}





\section{Deploying a TensorFlow Lite Model to Arduino}

n this tutorial, we will load the model in Arduino using the TensorFlow Lite library and use it to run inference to generate an approximation of a sinewave.

Note that this is a ridiculous, roundabout way to create a sinewave, but it offers a useful example to show how a nonlinear neural network model can be deployed to an embedded system.

You can watch the steps in this tutorial through video form \HREF{https://youtu.be/dU01M61RW8s}{here}.

\subsection{Model Description}

Our model is a 3-layer, fully connected neural network with a single, floating point input and a single, floating point output.

\begin{figure}
  \includegraphics[width=\textwidth]{TinyML/Short/NN}
  \caption{Neural network model}
\end{figure}

If you downloaded the .tflite file in the previous tutorial, you can use Netron to view the model in a slick, graphical interface. Run Netron and use it to open the .tflite file. You can click on the individual layers to get more details about them, such as input/output tensor shapes and data types.

\begin{figure}
  \includegraphics[width=\textwidth]{TinyML/Short/InspectModelWithNetron}
  \caption{Using Netron to view neural network model}
\end{figure}


As you can see, our model expects a tensor with 1 floating point element (a scalar value) as an input, and it outputs another scalar value. The input should be between $0$ and $2\pi$, and the output should be between $-1$ and $1$.

\subsection{Install TensorFlow Lite Arduino Library}

TensorFlow Lite has support for a few microcontroller boards, which are listed here. At the time this tutorial was released, only 8 microcontroller boards were supported. We will use the pre-compiled TensorFlow Lite library for TensorFlow Lite, but note that only the Nano 33 BLE Sense is supported (right now). 

Open your Arduino IDE (this tutorial was tested on v1.8.11). Go to Sketch $>$ Include Library $>$ Manage Librarie,$\ldots$ and search for ``TensorFlow''. Install the latest version of the \FILE{Arduino\_TensorFlowLite} library (1.15.0-ALPHA was tested for this tutorial). 

\begin{figure}
  \includegraphics[width=\textwidth]{TinyML/Short/ArduinoIDENano33}
  \caption{Install TensorFlow Lite library for Arduino}
\end{figure}



In Tools $>$ Board, select the Arduino Nano 33 BLE. Plug in your Nano 33 BLE Sense and select the associated Serial port in Tools $>$ Port.

\subsection{Test Inference}

Note: This code was originally developed by Pete Warden of the TensorFlow Lite team to demonstrate TensorFlow Lite capabilities on various microcontroller platforms. I have modified it here to reduce the number of dependent files in Arduino, which should hopefully make it easier to follow.

Copy the following into your Arduino sketch:



\medskip

 \textbf{Copy Code}
 
\begin{lstlisting} 
/**
* Test sinewave neural network model
* 
* Author: Pete Warden
* Modified by: Shawn Hymel
* Date: March 11, 2020
* 
* Copyright 2019 The TensorFlow Authors. All Rights Reserved.
*
* Licensed under the Apache License, Version 2.0 (the "License");
* you may not use this file except in compliance with the License.
* You may obtain a copy of the License at
*
*    http://www.apache.org/licenses/LICENSE-2.0
* 
* Unless required by applicable law or agreed to in writing, software
* distributed under the License is distributed on an "AS IS" BASIS,
* WITHOUT WARRANTIES OR CONDITIONS OF ANY KIND, either express or implied.
* See the License for the specific language governing permissions and
* limitations under the License.
*/

// Import TensorFlow stuff
#include "TensorFlowLite.h"
#include "tensorflow/lite/experimental/micro/kernels/micro_ops.h"
#include "tensorflow/lite/experimental/micro/micro_error_reporter.h"
#include "tensorflow/lite/experimental/micro/micro_interpreter.h"
#include "tensorflow/lite/experimental/micro/micro_mutable_op_resolver.h"
#include "tensorflow/lite/version.h"

// Our model
#include "sine_model.h"

// Figure out what's going on in our model
#define DEBUG 1

// Some settings
constexpr int led_pin = 2;
constexpr float pi = 3.14159265;                  // Some pi
constexpr float freq = 0.5;                       // Frequency (Hz) of sinewave
constexpr float period = (1 / freq) * (1000000);  // Period (microseconds)

// TFLite globals, used for compatibility with Arduino-style sketches
namespace {
    tflite::ErrorReporter* error_reporter = nullptr;
    const tflite::Model* model = nullptr;
    tflite::MicroInterpreter* interpreter = nullptr;
    TfLiteTensor* model_input = nullptr;
    TfLiteTensor* model_output = nullptr;
    
    // Create an area of memory to use for input, output, and other TensorFlow
    // arrays. You'll need to adjust this by combiling, running, and looking
    // for errors.
    constexpr int kTensorArenaSize = 5 * 1024;
    uint8_t tensor_arena[kTensorArenaSize];
} // namespace

void setup() {
    
    // Wait for Serial to connect
    #if DEBUG
    while(!Serial);
    #endif
    
    // Let's make an LED vary in brightness
    pinMode(led_pin, OUTPUT);
    
    // Set up logging (will report to Serial, even within TFLite functions)
    static tflite::MicroErrorReporter micro_error_reporter;
    error_reporter = &micro_error_reporter;
    
    // Map the model into a usable data structure
    model = tflite::GetModel(sine_model);
    if (model->version() != TFLITE_SCHEMA_VERSION) {
        error_reporter->Report("Model version does not match Schema");
        while(1);
    }
    
    // Pull in only needed operations (should match NN layers)
    // Available ops:
    //  https://github.com/tensorflow/tensorflow/blob/master/tensorflow/lite/micro/kernels/micro_ops.h
    static tflite::MicroMutableOpResolver micro_mutable_op_resolver;
    micro_mutable_op_resolver.AddBuiltin(
    tflite::BuiltinOperator_FULLY_CONNECTED,
    tflite::ops::micro::Register_FULLY_CONNECTED(),
    1, 3);
    
    // Build an interpreter to run the model
    static tflite::MicroInterpreter static_interpreter(
    model, micro_mutable_op_resolver, tensor_arena, kTensorArenaSize,
    error_reporter);
    interpreter = &static_interpreter;
    
    // Allocate memory from the tensor_arena for the model's tensors
    TfLiteStatus allocate_status = interpreter->AllocateTensors();
    if (allocate_status != kTfLiteOk) {
        error_reporter->Report("AllocateTensors() failed");
        while(1);
    }
    
    // Assign model input and output buffers (tensors) to pointers
    model_input = interpreter->input(0);
    model_output = interpreter->output(0);
    
    // Get information about the memory area to use for the model's input
    // Supported data types:
    // https://github.com/tensorflow/tensorflow/blob/master/tensorflow/lite/c/common.h#L226
    #if DEBUG
    Serial.print("Number of dimensions: ");
    Serial.println(model_input->dims->size);
    Serial.print("Dim 1 size: ");
    Serial.println(model_input->dims->data[0]);
    Serial.print("Dim 2 size: ");
    Serial.println(model_input->dims->data[1]);
    Serial.print("Input type: ");
    Serial.println(model_input->type);
    #endif
}

void loop() {
    
    #if DEBUG
    unsigned long start_timestamp = micros();
    #endif
    
    // Get current timestamp and modulo with period
    unsigned long timestamp = micros();
    timestamp = timestamp % (unsigned long)period;
    
    // Calculate x value to feed to the model
    float x_val = ((float)timestamp * 2 * pi) / period;
    
    // Copy value to input buffer (tensor)
    model_input->data.f[0] = x_val;
    
    // Run inference
    TfLiteStatus invoke_status = interpreter->Invoke();
    if (invoke_status != kTfLiteOk) {
        error_reporter->Report("Invoke failed on input: %f\n", x_val);
    }
    
    // Read predicted y value from output buffer (tensor)
    float y_val = model_output->data.f[0];
    
    // Translate to a PWM LED brightness
    int brightness = (int)(255 * y_val);
    analogWrite(led_pin, brightness);
    
    // Print value
    Serial.println(y_val);
    
    #if DEBUG
    Serial.print("Time for inference (us): ");
    Serial.println(micros() - start_timestamp);
    #endif
}
\end{lstlisting}



For this first test, we're simply assigning the value of pi (approximated to 3.14159265) to the input tensor:



\medskip 

\textbf{Copy Code}

\begin{lstlisting}
    model_input->data.f[0] = pi;
\end{lstlisting}



With this, we will run inference on just one number (over and over again, as it is in the loop function).

When you run the sketch, open the Serial monitor to see the output value. Also, note that because we have the \PYTHON{DEBUG} flag defined to \PYTHON{1}, we can get an estimate of how long inference takes.

\begin{figure}
  \includegraphics[width=\textwidth]{TinyML/Short/ModelOutput}
  \caption{Model output}
\end{figure}



Notice that the output of the model is always $0.03$. While $\sin(\pi)$ should be $0$, $0.03$ seems to be close enough. Additionally, notice that it takes about $1$ ms to perform inference. This number is based on the performance capabilities of the microcontroller and the size/complexity of the model. This is a relatively simple model, so expect inference to take longer as you start using more complex models.

\subsection{Examine the Code}

Let's take a quick look at what's going on in the code.

At the top, we define a number of pointers that we'll use throughout the rest of the sketch. Note that they're in an anonymous namespace, which probably isn't needed but follows the other TensorFlow Lite examples.



\medskip

 \textbf{Copy Code}
 
 \begin{lstlisting}
// TFLite globals, used for compatibility with Arduino-style sketches
namespace {
    tflite::ErrorReporter* error_reporter = nullptr;
    const tflite::Model* model = nullptr;
    tflite::MicroInterpreter* interpreter = nullptr;
    TfLiteTensor* model_input = nullptr;
    TfLiteTensor* model_output = nullptr;
    
    // Create an area of memory to use for input, output, and other TensorFlow
    // arrays. You'll need to adjust this by combiling, running, and looking
    // for errors.
    constexpr int kTensorArenaSize = 5 * 1024;
    uint8_t tensor_arena[kTensorArenaSize];
} // namespace
\end{lstlisting}



Note that we are setting aside a chunk of memory for an ``arena'' (essentially, a sandbox of RAM that TensorFlow Lite uses to perform calculations and store tensors). Unfortunately, we must predict the arena size. 5 kB seems to work for this model, but if you have problems during the ``allocate tensors'' step later, you should try increasing the arena size.

We set up logging, which, for Arduino, will output debugging information to the Serial port. Note that this will output information from within some of the TensorFlow Lite functions, which can be very helpful when something isn't working.



\medskip

\textbf{Copy Code}

\begin{lstlisting}
// Set up logging (will report to Serial, even within TFLite functions)
static tflite::MicroErrorReporter micro_error_reporter;
error_reporter = &micro_error_reporter;
\end{lstlisting}



To help reduce the required space, we only pull in the required TensorFlow Lite operations with the following code:



\medskip 

\textbf{Copy Code}


\begin{lstlisting}
// Pull in only needed operations (should match NN layers)
static tflite::MicroMutableOpResolver micro_mutable_op_resolver;
micro_mutable_op_resolver.AddBuiltin(
tflite::BuiltinOperator_FULLY_CONNECTED,
tflite::ops::micro::Register_FULLY_CONNECTED(),
1, 3);
\end{lstlisting}




You will need to pick out the ones that line up with the layers (and other operations) you defined when constructing the model. It can help to look at Netron to see the required operations.

In this next section, we build the interpreter using the parameters we just defined and then allocate the necessary memory using the arena. Then, we get handles to the input and output tensors. Note that these tensors are just one element each (scalars).



\medskip

\textbf{Copy Code}

\begin{lstlisting}
// Build an interpreter to run the model
static tflite::MicroInterpreter static_interpreter(
model, micro_mutable_op_resolver, tensor_arena, kTensorArenaSize,
error_reporter);
interpreter = &static_interpreter;

// Allocate memory from the tensor_arena for the model's tensors
TfLiteStatus allocate_status = interpreter->AllocateTensors();
if (allocate_status != kTfLiteOk) {
    error_reporter->Report("AllocateTensors() failed");
    while(1);
}

// Assign model input and output buffers (tensors) to pointers
model_input = interpreter->input(0);
model_output = interpreter->output(0);
\end{lstlisting}



If the \PYTHON{AllocateTensors()} function fails during runtime, you may want to try increasing the arena size defined earlier. There seems to be a bit of trial and error required to get the memory management right.

In loop, we copy our input values to the input tensor. Note the index of 0; if this was a multi-element tensor, we’d want to copy everything in with something like a for loop.



\medskip

 \textbf{Copy Code}
 
\begin{lstlisting}
// Copy value to input buffer (tensor)
model_input->data.f[0] = pi;
\end{lstlisting}



We tell the interpreter to run inference with the Invoke() function:



\medskip

 \textbf{Copy Code}
 
\begin{lstlisting}
// Run inference
TfLiteStatus invoke_status = interpreter->Invoke();
if (invoke_status != kTfLiteOk) {
    error_reporter->Report("Invoke failed on input: %f\n", x_val);
}
\end{lstlisting}



Note that at this time, Invoke is blocking, so we must wait while it performs its calculations.

When it’s done, we can access the output value in our \PYTHON{model\_output} handle:


\medskip

\textbf{Copy Code}

\begin{lstlisting}
float y_val = model_output->data.f[0];
\end{lstlisting}


Hopefully, this should give you an idea of the required TensorFlow operations that you need to run to get inference working on a microcontroller!

\subsection{Connect Inference to Hardware}


Connect an LED from pin 2 to GND (through a $100 \Omega$). 

At the top of the program, change the \PYTHON{DEBUG} flag to the following:



\medskip 

\textbf{Copy Code}

\begin{lstlisting}
#define DEBUG 0
\end{lstlisting}




Then, change the \PYTHON{loop()} function to the following:



\medskip

\textbf{Copy Code}


\begin{lstlisting}
void loop() {
    
    #if DEBUG
    unsigned long start_timestamp = micros();
    #endif
    
    // Get current timestamp and modulo with period
    unsigned long timestamp = micros();
    timestamp = timestamp % (unsigned long)period;
    
    // Calculate x value to feed to the model
    float x_val = ((float)timestamp * 2 * pi) / period;
    
    // Copy value to input buffer (tensor)
    model_input->data.f[0] = x_val;
    
    // Run inference
    TfLiteStatus invoke_status = interpreter->Invoke();
    if (invoke_status != kTfLiteOk) {
        error_reporter->Report("Invoke failed on input: %f\n", x_val);
    }
    
    // Read predicted y value from output buffer (tensor)
    float y_val = model_output->data.f[0];
    
    // Translate to a PWM LED brightness
    int brightness = (int)(255 * y_val);
    analogWrite(led_pin, brightness);
    
    // Print value
    Serial.println(y_val);
    
    #if DEBUG
    Serial.print("Time for inference (us): ");
    Serial.println(micros() - start_timestamp);
    #endif
}
\end{lstlisting}



Notice that we're now calculating the input to the sine prediction model based on the current timestamp. By doing this, we can effectively control the sinewave’s frequency using the freq variable at the top of the program.

Upload this code and open a Serial monitor. You should see the values of the model’s output fly by.


\begin{figure}
  \includegraphics[width=\textwidth]{TinyML/Short/SerialMonitorOutput}
    \caption{Serial monitor with sine output}
\end{figure}



This isn’t terribly useful, so go close it and select Tools > Serial Plotter. This should show you a slow-moving sinewave as a plot versus time.

\begin{figure}
  \includegraphics[width=\textwidth]{TinyML/Short/SerialPlotterSineWave}
  \caption{Serial plotter with sinewave}
\end{figure}



If you take a look at the LED connected to the Arduino board, you should see it increasing and decreasing in brightness in something that approximates a sinusoidal pattern.

\begin{figure}
  \includegraphics[width=\textwidth]{TinyML/Short/ArduinoNano33}
  \caption{LED slowly blinking in sinewave pattern}
\end{figure}


Try changing the frequency to something higher, say:



\medskip

\textbf{Copy Code}


\begin{lstlisting}
constexpr float freq = 100;
\end{lstlisting}




Upload this and look at the Serial Plotter. You should see a much better representation of our sinewave.

\begin{figure}
  \includegraphics[width=\textwidth]{TinyML/Short/SineWave}
  \caption{100 Hz sinewave on Serial Plotter}
\end{figure}





 \subsection{Going Further}

Hopefully, this has helped you get started using Colab to create a simple TensorFlow Lite model with the intention of deploying it to a microcontroller! On the next tutorial, we will run the TensorFlow Lite inference engine on an Arduino and use our model to predict sine function values.

See these articles to learn more about TensorFlow Lite for Microcontrollers:

\begin{itemize}
    \item \URL{https://www.tensorflow.org/lite/microcontrollers}
    \item \URL{https://www.tensorflow.org/lite/microcontrollers/get_started}
    \item \URL{https://www.digikey.com/en/maker/projects/tensorflow-lite-for-microcontrollers-kit-quickstart/1b372b69e44f4d988b5363741a61882d}
\end{itemize}

\subsection{Key Parts and Components}

\begin{itemize}
    \item Arduino ARDUINO NANO 33 BLE SENSE, ABX00031, \EUR{45,50}
\end{itemize}

